\chapter{Revisão Prática da Parte III: SQL Aplicada}
\label{cap:revisao_pratica_sql_aplicada}

A Parte III conduziu o leitor do aprofundamento em SQL até temas de administração inicial de banco de dados. Ao longo das Aulas 10 a 17, foram estudados os principais grupos da linguagem SQL, a definição e modificação de estruturas, a manipulação de dados, a construção de consultas analíticas, o uso de views, o controle de acesso, os índices, as transações e as funções.

Este capítulo final tem como objetivo consolidar esse percurso. Em vez de introduzir um novo bloco teórico isolado, a proposta é revisar, integrar e aplicar os conteúdos da Parte III em um encadeamento prático. Assim, o leitor pode perceber como os diversos recursos estudados se articulam em um cenário mais próximo do uso real de um banco de dados relacional.

\section{Panorama de fechamento da Parte III}

Ao final desta parte, o estudante deve ser capaz de:

\begin{itemize}
    \item compreender a SQL como linguagem de implementação e uso de bancos relacionais;
    \item criar estruturas com DDL;
    \item manipular dados com DML;
    \item elaborar consultas mais sofisticadas com agregações, junções e subconsultas;
    \item refinar consultas com ordenação, paginação, remoção de duplicidades e funções de janela;
    \item utilizar \texttt{VIEW}s como mecanismo de abstração;
    \item aplicar noções de segurança com roles e privilégios;
    \item entender o papel de índices no desempenho;
    \item usar transações para proteger operações compostas;
    \item encapsular lógicas recorrentes com funções.
\end{itemize}

Essa lista mostra que a Parte III não se limita a “comandos de SQL”, mas constitui um bloco de transição entre o aprendizado da linguagem e sua aplicação profissional.

\section{Mapa conceitual resumido}

Uma forma útil de visualizar a Parte III é organizar seus tópicos em quatro eixos:

\subsection{Eixo estrutural}

Relaciona-se à criação e modificação da estrutura do banco:
\begin{itemize}
    \item \texttt{CREATE}
    \item \texttt{ALTER}
    \item \texttt{DROP}
    \item tipos de dados
    \item restrições
    \item integridade referencial
\end{itemize}

\subsection{Eixo operacional}

Relaciona-se à entrada e manutenção dos dados:
\begin{itemize}
    \item \texttt{INSERT}
    \item \texttt{SELECT}
    \item \texttt{UPDATE}
    \item \texttt{DELETE}
\end{itemize}

\subsection{Eixo analítico}

Relaciona-se à produção de respostas mais sofisticadas:
\begin{itemize}
    \item agregações
    \item \texttt{GROUP BY}
    \item \texttt{HAVING}
    \item subconsultas
    \item \texttt{JOIN}s
    \item \texttt{ORDER BY}
    \item \texttt{LIMIT}
    \item \texttt{DISTINCT}
    \item \texttt{UNION}
    \item funções de janela
    \item \texttt{VIEW}
\end{itemize}

\subsection{Eixo administrativo}

Relaciona-se à manutenção do ambiente e ao comportamento do banco:
\begin{itemize}
    \item roles e privilégios
    \item \texttt{GRANT} e \texttt{REVOKE}
    \item índices
    \item transações
    \item funções
\end{itemize}

\section{Estudo de caso integrador}

Para consolidar a Parte III, será utilizado um pequeno sistema acadêmico com as seguintes tabelas:

\begin{itemize}
    \item \texttt{cursos}
    \item \texttt{alunos}
    \item \texttt{disciplinas}
    \item \texttt{notas}
    \item \texttt{turma}
    \item \texttt{matricula}
\end{itemize}

A ideia não é esgotar todas as possibilidades do sistema, mas mostrar como os diversos recursos da Parte III se combinam em uma sequência coerente.

\section{Passo 1 -- Estrutura inicial}

\subsection{Criação de cursos}

\begin{sqlcode}
CREATE TABLE cursos (
    curso_id SERIAL PRIMARY KEY,
    nome VARCHAR(100) NOT NULL,
    coordenador VARCHAR(100)
);
\end{sqlcode}

\subsection{Criação de alunos}

\begin{sqlcode}
CREATE TABLE alunos (
    aluno_id SERIAL PRIMARY KEY,
    nome VARCHAR(100) NOT NULL,
    idade INTEGER,
    curso_id INTEGER,
    status VARCHAR(20) DEFAULT 'ATIVO',
    FOREIGN KEY (curso_id) REFERENCES cursos(curso_id)
);
\end{sqlcode}

\subsection{Criação de disciplinas}

\begin{sqlcode}
CREATE TABLE disciplinas (
    disciplina_id SERIAL PRIMARY KEY,
    nome VARCHAR(100) NOT NULL
);
\end{sqlcode}

\subsection{Criação de notas}

\begin{sqlcode}
CREATE TABLE notas (
    nota_id SERIAL PRIMARY KEY,
    aluno_id INTEGER NOT NULL,
    disciplina_id INTEGER NOT NULL,
    nota NUMERIC(5,2) NOT NULL,
    FOREIGN KEY (aluno_id) REFERENCES alunos(aluno_id),
    FOREIGN KEY (disciplina_id) REFERENCES disciplinas(disciplina_id)
);
\end{sqlcode}

\subsection{Criação de turma}

\begin{sqlcode}
CREATE TABLE turma (
    turma_id SERIAL PRIMARY KEY,
    nome VARCHAR(50) NOT NULL,
    vagas_disponiveis INTEGER NOT NULL DEFAULT 30
);
\end{sqlcode}

\subsection{Criação de matrícula}

\begin{sqlcode}
CREATE TABLE matricula (
    matricula_id SERIAL PRIMARY KEY,
    aluno_id INTEGER NOT NULL,
    turma_id INTEGER NOT NULL,
    FOREIGN KEY (aluno_id) REFERENCES alunos(aluno_id),
    FOREIGN KEY (turma_id) REFERENCES turma(turma_id)
);
\end{sqlcode}

\section{Passo 2 -- Carga inicial de dados}

\subsection{Inserindo cursos}

\begin{sqlcode}
INSERT INTO cursos (nome, coordenador) VALUES
('Engenharia de Software', 'Marcos Silva'),
('Sistemas de Informação', 'Luciana Costa'),
('Ciência da Computação', 'Carlos Mendes');
\end{sqlcode}

\subsection{Inserindo alunos}

\begin{sqlcode}
INSERT INTO alunos (nome, idade, curso_id) VALUES
('Aline Castro', 20, 1),
('Rafael Gomes', 21, 2),
('Bianca Luz', 22, 3),
('Marcelo Reis', 19, 1),
('Paula Souza', 23, 2);
\end{sqlcode}

\subsection{Inserindo disciplinas}

\begin{sqlcode}
INSERT INTO disciplinas (nome) VALUES
('Banco de Dados'),
('Algoritmos'),
('Estruturas de Dados');
\end{sqlcode}

\subsection{Inserindo notas}

\begin{sqlcode}
INSERT INTO notas (aluno_id, disciplina_id, nota) VALUES
(1, 1, 7.5),
(2, 2, 8.0),
(3, 1, 6.5),
(4, 3, 9.0),
(5, 1, 8.8);
\end{sqlcode}

\subsection{Inserindo turmas}

\begin{sqlcode}
INSERT INTO turma (nome, vagas_disponiveis) VALUES
('ESW-2025-1A', 30),
('SIN-2025-1A', 25),
('CCO-2025-1A', 20);
\end{sqlcode}

\section{Passo 3 -- Consultas operacionais}

\subsection{Listando alunos}

\begin{sqlcode}
SELECT *
FROM alunos;
\end{sqlcode}

\subsection{Listando alunos ativos}

\begin{sqlcode}
SELECT nome, idade
FROM alunos
WHERE status = 'ATIVO';
\end{sqlcode}

\subsection{Atualizando status}

\begin{sqlcode}
UPDATE alunos
SET status = 'INATIVO'
WHERE aluno_id = 5;
\end{sqlcode}

\subsection{Removendo nota lançada indevidamente}

\begin{sqlcode}
DELETE FROM notas
WHERE nota_id = 3;
\end{sqlcode}

\textbf{Comentário:}  
Essas operações retomam diretamente a DML estudada nas Aulas 13 e 14, mostrando a manutenção cotidiana dos registros.

\section{Passo 4 -- Consultas analíticas}

\subsection{Quantidade de alunos por curso}

\begin{sqlcode}
SELECT
    curso_id,
    COUNT(*) AS total_alunos
FROM alunos
GROUP BY curso_id;
\end{sqlcode}

\subsection{Cursos com pelo menos dois alunos}

\begin{sqlcode}
SELECT
    curso_id,
    COUNT(*) AS total_alunos
FROM alunos
GROUP BY curso_id
HAVING COUNT(*) >= 2;
\end{sqlcode}

\subsection{Alunos acima da média de idade}

\begin{sqlcode}
SELECT nome, idade
FROM alunos
WHERE idade > (
    SELECT AVG(idade)
    FROM alunos
);
\end{sqlcode}

\subsection{Relatório com aluno, disciplina e nota}

\begin{sqlcode}
SELECT a.nome AS aluno,
       d.nome AS disciplina,
       n.nota
FROM notas n
INNER JOIN alunos a
    ON n.aluno_id = a.aluno_id
INNER JOIN disciplinas d
    ON n.disciplina_id = d.disciplina_id;
\end{sqlcode}

\subsection{Cursos sem alunos}

\begin{sqlcode}
SELECT c.nome
FROM cursos c
LEFT JOIN alunos a
    ON c.curso_id = a.curso_id
WHERE a.aluno_id IS NULL;
\end{sqlcode}

\textbf{Comentário:}  
Aqui aparecem, integradamente, agregações, subconsultas e junções, retomando o núcleo analítico das Aulas 14 e 15.

\section{Passo 5 -- Refinamento das consultas}

\subsection{Ordenação por curso e nome}

\begin{sqlcode}
SELECT nome, curso_id
FROM alunos
ORDER BY curso_id, nome;
\end{sqlcode}

\subsection{Paginação}

\begin{sqlcode}
SELECT nome, curso_id
FROM alunos
ORDER BY nome
LIMIT 3 OFFSET 0;
\end{sqlcode}

\subsection{Cursos distintos}

\begin{sqlcode}
SELECT DISTINCT curso_id
FROM alunos
ORDER BY curso_id;
\end{sqlcode}

\subsection{União entre alunos e coordenadores}

\begin{sqlcode}
SELECT nome, 'ALUNO' AS origem
FROM alunos
UNION ALL
SELECT coordenador, 'COORDENADOR' AS origem
FROM cursos;
\end{sqlcode}

\subsection{Média de idade do curso em cada linha}

\begin{sqlcode}
SELECT
    nome,
    curso_id,
    idade,
    AVG(idade) OVER (PARTITION BY curso_id) AS media_curso
FROM alunos;
\end{sqlcode}

\section{Passo 6 -- View para relatório acadêmico}

\begin{sqlcode}
CREATE VIEW vw_notas_alunos AS
SELECT
    a.nome AS aluno,
    d.nome AS disciplina,
    n.nota
FROM notas n
INNER JOIN alunos a
    ON n.aluno_id = a.aluno_id
INNER JOIN disciplinas d
    ON n.disciplina_id = d.disciplina_id;
\end{sqlcode}

\begin{sqlcode}
SELECT *
FROM vw_notas_alunos
ORDER BY aluno, disciplina;
\end{sqlcode}

\textbf{Comentário:}  
A view centraliza uma consulta que poderia ser usada repetidamente em relatórios e interfaces do sistema.

\section{Passo 7 -- Controle de acesso}

\subsection{Criando perfil de leitura}

\begin{sqlcode}
CREATE ROLE leitura_academica;
GRANT SELECT ON alunos TO leitura_academica;
GRANT SELECT ON cursos TO leitura_academica;
GRANT SELECT ON vw_notas_alunos TO leitura_academica;
\end{sqlcode}

\subsection{Criando usuário e vinculando perfil}

\begin{sqlcode}
CREATE ROLE usuario_consulta LOGIN PASSWORD 'consulta123';
GRANT leitura_academica TO usuario_consulta;
\end{sqlcode}

\textbf{Comentário:}  
Esse trecho recupera a Aula 16, mostrando que a disponibilização de dados também deve ser mediada por política de acesso.

\section{Passo 8 -- Índices}

\subsection{Índice para busca por nome do aluno}

\begin{sqlcode}
CREATE INDEX idx_alunos_nome
ON alunos(nome);
\end{sqlcode}

\subsection{Índice multicoluna para matrícula}

\begin{sqlcode}
CREATE INDEX idx_matricula_aluno_turma
ON matricula(aluno_id, turma_id);
\end{sqlcode}

\textbf{Comentário:}  
Esses índices são justificados por padrões de consulta plausíveis no sistema.

\section{Passo 9 -- Transação de matrícula}

\begin{sqlcode}
BEGIN;

INSERT INTO matricula (aluno_id, turma_id)
VALUES (1, 1);

UPDATE turma
SET vagas_disponiveis = vagas_disponiveis - 1
WHERE turma_id = 1;

COMMIT;
\end{sqlcode}

\textbf{Comentário:}  
A transação protege uma operação composta. Sem ela, o banco poderia registrar a matrícula sem atualizar corretamente a vaga disponível, ou vice-versa.

\section{Passo 10 -- Função de apoio}

\begin{sqlcode}
CREATE OR REPLACE FUNCTION total_matriculas_aluno(p_aluno_id INTEGER)
RETURNS INTEGER
LANGUAGE plpgsql
AS $$
DECLARE
    v_total INTEGER;
BEGIN
    SELECT COUNT(*)
    INTO v_total
    FROM matricula
    WHERE aluno_id = p_aluno_id;

    RETURN v_total;
END;
$$;
\end{sqlcode}

\begin{sqlcode}
SELECT total_matriculas_aluno(1);
\end{sqlcode}

\textbf{Comentário:}  
Essa função mostra como o banco pode oferecer operações reutilizáveis e não apenas armazenar dados brutos.

\section{Erros clássicos de integração}

Ao revisar toda a Parte III, alguns erros aparecem como especialmente recorrentes:

\subsection{Erro 1 -- Esquecer a cláusula WHERE em UPDATE ou DELETE}

Isso pode alterar ou remover mais linhas do que o pretendido.

\subsection{Erro 2 -- Esquecer a condição de junção}

Sem a condição adequada, o resultado pode virar um produto cartesiano.

\subsection{Erro 3 -- Confundir WHERE com HAVING}

\texttt{WHERE} filtra linhas; \texttt{HAVING} filtra grupos.

\subsection{Erro 4 -- Criar privilégio amplo demais}

Conceder \texttt{ALL PRIVILEGES} ou \texttt{SUPERUSER} sem necessidade compromete segurança.

\subsection{Erro 5 -- Criar índice sem necessidade real}

Índices devem refletir padrões concretos de uso, não apenas intuição.

\subsection{Erro 6 -- Não usar transação em operação composta}

Sem transação, o banco pode ficar inconsistente se apenas parte do processo for concluída.

\section{Boas práticas finais da Parte III}

Ao concluir esta parte do livro, algumas boas práticas merecem destaque:

\begin{itemize}
    \item validar operações destrutivas com consultas prévias;
    \item preferir clareza e legibilidade na escrita SQL;
    \item usar aliases consistentes;
    \item separar estrutura, manipulação, análise e segurança mentalmente;
    \item tratar o banco como ambiente ativo de regras e não apenas como depósito de dados;
    \item sempre pensar no impacto de desempenho, consistência e acesso.
\end{itemize}

\begin{quadro_dica}
Uma maneira madura de pensar SQL é enxergar cada comando em três planos ao mesmo tempo: o que ele faz, que risco ele envolve e que papel ele desempenha no sistema como um todo.
\end{quadro_dica}

\section{Exercícios de revisão}

\subsection{Bloco 1 -- Estrutura e manipulação}

\begin{enumerate}
    \item Crie uma tabela de exemplo com chave primária e chave estrangeira.
    \item Insira ao menos cinco registros.
    \item Atualize dois deles.
    \item Remova um registro de forma controlada.
\end{enumerate}

\subsection{Bloco 2 -- Consulta analítica}

\begin{enumerate}
    \item Conte a quantidade de registros por grupo.
    \item Calcule uma média por grupo.
    \item Filtre grupos por meio de \texttt{HAVING}.
    \item Elabore uma subconsulta com \texttt{EXISTS}.
    \item Escreva uma consulta com \texttt{INNER JOIN} e outra com \texttt{LEFT JOIN}.
\end{enumerate}

\subsection{Bloco 3 -- Refinamento e views}

\begin{enumerate}
    \item Faça uma consulta com \texttt{ORDER BY}, \texttt{LIMIT} e \texttt{OFFSET}.
    \item Liste valores distintos com \texttt{DISTINCT}.
    \item Use \texttt{UNION ALL} para combinar dois conjuntos.
    \item Crie uma função de janela simples.
    \item Crie uma \texttt{VIEW} para relatório.
\end{enumerate}

\subsection{Bloco 4 -- Administração}

\begin{enumerate}
    \item Crie uma role de leitura.
    \item Conceda privilégios coerentes.
    \item Crie um índice justificado.
    \item Escreva uma transação composta.
    \item Escreva uma função reutilizável.
\end{enumerate}

\section{Exercício final integrador}

Projete e implemente um pequeno sistema de \textbf{biblioteca universitária} com as tabelas:
\begin{itemize}
    \item \texttt{livro}
    \item \texttt{autor}
    \item \texttt{categoria}
    \item \texttt{leitor}
    \item \texttt{emprestimo}
\end{itemize}

Depois:

\begin{enumerate}
    \item insira dados de teste;
    \item escreva consultas operacionais;
    \item produza consultas analíticas com agrupamento;
    \item crie ao menos uma \texttt{VIEW};
    \item defina uma role de leitura;
    \item crie um índice justificado;
    \item escreva uma transação de empréstimo;
    \item implemente uma função de apoio;
    \item explique por que cada recurso foi utilizado.
\end{enumerate}

\section{Fechamento da Parte III}

A Parte III teve como eixo central o aprofundamento da SQL. Seu percurso começou com a consolidação da linguagem como ferramenta de implementação e avançou para o uso mais profissional dos bancos de dados relacionais. Ao final desta parte, o leitor não apenas conhece comandos isolados, mas também compreende como eles se articulam em estruturas, operações, relatórios, controles de acesso, decisões de desempenho e rotinas reutilizáveis.

Dessa forma, a SQL deixa de ser percebida apenas como linguagem de consulta e passa a ser entendida como instrumento completo de trabalho com bancos de dados: definição, manipulação, análise, proteção e administração inicial.

\section{Síntese do capítulo}

Este capítulo final consolidou as Aulas 10 a 17 por meio de uma revisão prática integrada. Foram retomados os principais conceitos da Parte III e apresentados em um fluxo único de trabalho, abrangendo DDL, DML, consultas avançadas, views, roles, índices, transações e funções. Com isso, encerra-se a Parte III com uma visão unificada da SQL aplicada, oferecendo ao leitor uma base mais robusta para estudos posteriores e para atuação prática em banco de dados.